% lecture03.tex — STAT 212, Lecture 3: Radon–Nikodym Theorem Proof, Martingales

\section*{Agenda}
\begin{enumerate}
    \item Radon--Nikodym Theorem
    \item Martingales
\end{enumerate}

\noindent\textbf{Announcement.} Lectures in G125 from now on.

%% ========================================================================
\section{Radon--Nikodym Theorem (continued)}
%% ========================================================================

\subsection{Recap}

So far: $(\Omega, \F, \Prob)$ is a measure (probability) space.
$\mu, \nu$ are finite measures on $(\Omega, \F)$, with $\nu(S) \leq \mu(S)$ for all $S \in \F$.

\medskip
\noindent\textbf{To prove:} There exists $f : \Omega \to [0,1]$ such that
\[
    \nu(S) = \int_S f \dmu.
\]

\noindent\textbf{Have done:} Define
\[
    \cH = \left\{ f : \Omega \to [0,1] \text{ measurable} \;\Big|\; \int_S f \dmu \leq \nu(S) \;\; \forall\, S \in \F \right\}, \qquad
    M = \sup_{f \in \cH} \int f \dmu.
\]

Then there exists $f_* \in \cH$ such that $\int f_* \dmu = M$.
(Recall: $f_*$ was constructed as the pointwise limit of $\bar{f}_n$, where $\bar{f}_n = \max(f_1, \ldots, f_n)$ for a maximizing sequence $\{f_n\}$ in $\cH$.)

\subsection{Proof that \texorpdfstring{$f_*$}{f*} achieves \texorpdfstring{$\nu$}{nu}}

\begin{claim}{$f_*$ satisfies $\nu(S) = \int_S f_* \dmu$}{rn-claim}
    For every $S \in \F$,
    \[
        \nu(S) = \int_S f_* \dmu.
    \]
\end{claim}

\begin{proof}
Suppose not. Then there exists $\bar{S} \in \F$ such that
\[
    \nu(\bar{S}) > \int_{\bar{S}} f_* \dmu.
\]

\noindent\textbf{Strategy:} Use $\bar{S}$ to construct $\hat{S} \in \F$ with $\mu(\hat{S}) > 0$ such that $f_* + \varepsilon \ind_{\hat{S}} \in \cH$.
But then
\[
    \int (f_* + \varepsilon \ind_{\hat{S}}) \dmu > M,
\]
which is a contradiction.

\medskip
\noindent\textbf{Construction strategy:} Iterative construction, carving out bad sets at each step.

\bigskip
\noindent\textbf{Step 1: Observations.}

\begin{enumerate}[label=(\roman*)]
    \item \emph{Observation on $\bar{S}$:}
    Since $\nu \leq \mu$ (by definition) and $f_* \in \cH$,
    \[
        \nu\bigl(S \cap \{f_* = 1\}\bigr) \leq \mu\bigl(S \cap \{f_* = 1\}\bigr), \quad \text{(by definition, } \nu \leq \mu\text{)},
    \]
    and
    \[
        \nu\bigl(S \cap \{f_* = 1\}\bigr) \geq \int_{S \cap \{f_* = 1\}} f_* \dmu = \mu\bigl(S \cap \{f_* = 1\}\bigr).
    \]
    So the deficit $\nu(\bar{S}) - \int_{\bar{S}} f_* \dmu$ is concentrated on $\bar{S} \cap \{f_* < 1\}$.

    \item \emph{Shrinking to $S_0$:}
    Start with $S_0 = \bar{S} \setminus \{f_* = 1\}$.
    Let $F_n = \{f_* \leq 1 - \tfrac{1}{n}\}$. Then $F_n \uparrow \{f_* < 1\}$.

    Since $\int_{\bar{S}} f_* \dmu < \nu(\bar{S})$, there exists $n_0$ large enough such that
    \[
        \int_{S \cap F_n} f_* \dmu < \nu(S \cap F_n), \qquad \forall\, n \geq n_0.
    \]

    So there exists $n$ large enough (set $\varepsilon = \varepsilon(n)$) such that:
    \begin{enumerate}[label=(\alph*)]
        \item $f_* + \varepsilon \ind_{S \cap F_n} \leq 1$ on $S \cap F_n$ (since $f_* \leq 1 - \tfrac{1}{n}$ on $F_n$, choosing $\varepsilon \leq \tfrac{1}{n}$),
        \item $\int_{S \cap F_n} (f_* + \varepsilon \ind_{S \cap F_n}) \dmu < \nu(S \cap F_n)$.
    \end{enumerate}

    Setting $\bar{S} = S \cap F_n$, the function $f_* + \varepsilon \ind_{\bar{S}}$ seems like a good start.

    \textbf{Issue:} Is $f_* + \varepsilon \ind_{\bar{S}} \in \cH$?
    We need $\int_T (f_* + \varepsilon \ind_{\bar{S}}) \dmu \leq \nu(T)$ for all $T \in \F$, not just $T = \bar{S}$.
\end{enumerate}

\bigskip
\noindent\textbf{Step 2: Shrinking $\bar{S}$ to $\hat{S}$.}

\emph{Idea:} Need to ``shrink'' $\bar{S}$ to $\hat{S}$ so that $f_* + \varepsilon \ind_{\hat{S}} \in \cH$, i.e.,
\[
    \int_S (f_* + \varepsilon \ind_{\hat{S}}) \dmu \leq \nu(S), \qquad \forall\, S \in \F.
\]

\emph{Observation:} If $T \subseteq \bar{S}^c$, then
\[
    \int_T (f_* + \varepsilon \ind_{\bar{S}}) \dmu \leq \nu(T),
\]
since the indicator does not matter there. So we need to worry only about $T \subseteq \bar{S}$.

\bigskip
\noindent\textbf{Step 3: The Deficit.}

\begin{definition}{Deficit}{deficit}
    For $A \in \F$, define
    \[
        \mathrm{Def}_\varepsilon(A) = \nu(A) - \int_A (f_* + \varepsilon \ind_A) \dmu.
    \]
    Note that $\mathrm{Def}_\varepsilon$ is additive (as a function of $A$, for sets contained in $\bar{S}$).
\end{definition}

By construction, $\mathrm{Def}_\varepsilon(\bar{S}) > 0$.

\medskip
\emph{Idea:} If $T \subseteq \bar{S}$ violates the deficit, i.e., $\mathrm{Def}_\varepsilon(T) < 0$, then
\[
    \mathrm{Def}_\varepsilon(\bar{S} \setminus T) = \mathrm{Def}_\varepsilon(\bar{S}) - \mathrm{Def}_\varepsilon(T) > 0.
\]

So: we will progressively rule out the bad sets.

\bigskip
\noindent\textbf{Step 4: Progressive elimination of bad sets.}

Define
\[
    a_k = \inf_{S_k} \mathrm{Def}_\varepsilon(S_k), \qquad S_k \subseteq \bar{S},
\]
where $S_k$ is chosen to be disjoint from the union of $S_1, \ldots, S_{k-1}$.

Pick $S_k$ such that $\mathrm{Def}_\varepsilon(S_k) \leq a_k + \tfrac{1}{k}$.
(Observe: $a_k \leq 0$ since $\emptyset$ is always allowed.)

Define:
\[
    \hat{S} = \bar{S} \setminus \bigcup_{k \geq 1} S_k.
\]

Then $\mathrm{Def}_\varepsilon(\hat{S}) > 0$, which implies $\nu(\hat{S}) > 0$, hence $\mu(\hat{S}) > 0$.

\bigskip
\noindent\textbf{Step 5: Verifying $f_* + \varepsilon \ind_{\hat{S}} \in \cH$.}

\begin{claim}{$f_* + \varepsilon \ind_{\hat{S}} \in \cH$}{fstar-in-H}
    The function $f_* + \varepsilon \ind_{\hat{S}}$ belongs to $\cH$.
\end{claim}

Suppose not. Then there exists $T \subseteq \hat{S}$ such that
\[
    \int_T (f_* + \varepsilon \ind_{\hat{S}}) \dmu > \nu(T) + \delta
\]
for some $\delta > 0$, i.e., $\mathrm{Def}_\varepsilon(T) < -\delta$.

Since $T \subseteq \hat{S}$, we have $T \cap S_k = \emptyset$ for all $k \geq 1$. So $T$ could have been used as $S_k$ for $k > 1/\delta$.

But then for $k$ large enough,
\[
    S_k \cup T \implies \mathrm{Def}_\varepsilon(T \cup S_k) \leq \mathrm{Def}_\varepsilon(S_k) - \delta \leq a_k + \frac{1}{k} - \delta < a_k,
\]
where the last inequality holds for $k > 1/\delta$. This is a contradiction with the definition of $a_k$ as the infimum!

\medskip
Therefore $f_* + \varepsilon \ind_{\hat{S}} \in \cH$, and
\[
    \int (f_* + \varepsilon \ind_{\hat{S}}) \dmu > M,
\]
contradicting the definition of $M$. This completes the proof.
\end{proof}

%% ========================================================================
\section{Conditional Probability and Conditional Distribution}
%% ========================================================================

So far: $(\Omega, \F, \Prob)$ probability space, $\G \subseteq \F$ a sub-$\sigma$-field.
If $\E\abs{X} < \infty$, we can define $\E[X \mid \G]$.

\medskip
\noindent\textbf{Q.} What about conditional probability? We can use $\E[\ind_A \mid \G]$.

\medskip
\noindent\textbf{Q.} What about conditional distribution? Same idea, but need to worry a bit more about measure-zero sets. (More in Sections.)

%% ========================================================================
\section{Martingales}
%% ========================================================================

\subsection{Setup}

$(\Omega, \F, \Prob)$ probability space. $\{\F_n : n \geq 1\}$ is a \emph{filtration}:
\begin{itemize}
    \item $\F_n$ is a $\sigma$-algebra, $\F_n \subseteq \F$,
    \item $\F_n \subseteq \F_{n+1}$.
\end{itemize}

$\{X_n : n \geq 1\}$ is \emph{adapted} to $\F_n$, i.e.,
\[
    \{X_n \in A\} \in \F_n.
\]

\begin{definition}{Martingale, Sub-martingale, Super-martingale}{mart-def}
    We say $\{(X_n, \F_n) : n \geq 1\}$ is a
    \begin{enumerate}[label=(\roman*)]
        \item \textbf{martingale} if $\E[X_{n+1} \mid \F_n] = X_n$,
        \item \textbf{sub-martingale} if $\E[X_{n+1} \mid \F_n] \geq X_n$,
        \item \textbf{super-martingale} if $\E[X_{n+1} \mid \F_n] \leq X_n$.
    \end{enumerate}
\end{definition}

\begin{example}{Random walk martingale}{rw-mart}
    Let $\varepsilon_i$ be i.i.d.\ with $\E[\varepsilon_i] = 0$ and $\E[\varepsilon_i^2] = 1$. Define
    \[
        S_k = \varepsilon_1 + \varepsilon_2 + \cdots + \varepsilon_k.
    \]
    Then $\{S_k\}$ is a martingale with respect to $\F_n = \sigma(\{S_0, S_1, \ldots, S_n\})$.
\end{example}

\subsection{Submartingale Convergence Theorem}

\begin{theorem}{Submartingale Convergence Theorem}{submart-conv}
    If $\{(X_n, \F_n) : n \geq 1\}$ is a sub-martingale with
    \[
        \sup_{n \geq 1} \E[X_n^+] < \infty,
    \]
    then $X_n \asto X_\infty$.
\end{theorem}

\noindent\textbf{Q.} What can we say about the limit? For example, what about $\E[X_\infty]$? $\E[X_\infty^m]$?

\medskip
\noindent\textbf{Naive hope:} If $\{X_n : n \geq 1\}$ is a martingale, then $\E[X_n] = \E[X_1]$ for all $n \geq 1$. Maybe $\E[X_n] = \E[X_\infty]$?

\medskip
Unfortunately no! Will need additional conditions. Next time.
